\documentclass[10pt,a4paper]{article}

% Encodage, langue et typographie
\usepackage[T1]{fontenc}
\usepackage[utf8]{inputenc}
\usepackage[french]{babel}
\usepackage{lmodern}
\usepackage{microtype}
\usepackage{ragged2e}

% Mise en page
\usepackage[a4paper,left=1.15cm,right=1.15cm,top=1.05cm,bottom=1.05cm]{geometry}
\usepackage{xcolor}
\usepackage{tikz}
\usepackage[most]{tcolorbox}
\usepackage{enumitem}
\usepackage{fontawesome5}
\usepackage{tabularx}
\usepackage[hidelinks]{hyperref}

% Palette sobre inspirée des CV français contemporains
\definecolor{navy}{HTML}{173B57}
\definecolor{blue}{HTML}{247BA0}
\definecolor{ice}{HTML}{EDF5F7}
\definecolor{ink}{HTML}{26343D}
\definecolor{muted}{HTML}{687780}
\definecolor{line}{HTML}{C9D8DE}

\hypersetup{
  colorlinks=true,
  urlcolor=navy,
  linkcolor=navy,
  pdfauthor={Camille Laurent},
  pdftitle={CV de Camille Laurent},
  pdfsubject={Cheffe de projet digital}
}

\pagestyle{empty}
\setlength{\parindent}{0pt}
\setlength{\parskip}{0pt}
\setlength{\fboxsep}{0pt}
\renewcommand{\familydefault}{\sfdefault}
\color{ink}

\setlist[itemize]{
  label={\color{blue}\raisebox{0.15ex}{\scriptsize\faCircle}},
  leftmargin=12pt,
  itemsep=1.5pt,
  topsep=3pt,
  parsep=0pt,
  partopsep=0pt
}

% Titres de rubriques
\newcommand{\mainsection}[1]{%
  \vspace{5.5pt}%
  {\large\bfseries\color{navy}\MakeUppercase{#1}}\par
  \vspace{1.5pt}%
  {\color{blue}\rule{\linewidth}{1.15pt}}\par
  \vspace{4pt}%
}

\newcommand{\sidesection}[1]{%
  \vspace{7pt}%
  {\normalsize\bfseries\color{navy}\MakeUppercase{#1}}\par
  \vspace{1.5pt}%
  {\color{blue}\rule{\linewidth}{0.8pt}}\par
  \vspace{4pt}%
}

% Entrée d'expérience : poste, dates, entreprise, ville
\newcommand{\experience}[4]{%
  \textbf{\color{ink}#1}\hfill
  {\footnotesize\bfseries\color{muted}#2}\par
  {\small\bfseries\color{blue}#3}%
  \hfill{\footnotesize\color{muted}\faMapMarker*\ #4}\par
}

% Entrée de formation : diplôme, dates, établissement, précision
\newcommand{\education}[4]{%
  \textbf{\color{ink}#1}\hfill
  {\footnotesize\bfseries\color{muted}#2}\par
  {\small\color{blue}#3}\par
  {\footnotesize\color{muted}#4}\par
}

% Compétence avec niveau visuel sur cinq points
\newcommand{\skill}[2]{%
  \begin{tabularx}{\linewidth}{@{}>{\raggedright\arraybackslash\bfseries\small}X r@{}}
    #1 &
    \foreach \n in {1,...,5}{%
      \ifnum\n>#2
        \textcolor{line}{\scriptsize\faCircle}%
      \else
        \textcolor{blue}{\scriptsize\faCircle}%
      \fi
      \hspace{1.2pt}%
    }
  \end{tabularx}\vspace{1pt}%
}

\newcommand{\contact}[2]{%
  \makebox[14pt][c]{\color{blue}\faIcon{#1}}%
  \hspace{3pt}#2%
}

\begin{document}

% En-tête
\begin{tcolorbox}[
  colback=navy,
  colframe=navy,
  boxrule=0pt,
  arc=2mm,
  left=7mm,
  right=7mm,
  top=5.5mm,
  bottom=5.5mm,
  enhanced
]
  \begin{minipage}[c]{0.16\linewidth}
    \centering
    \begin{tikzpicture}
      \fill[white] (0,0) circle (1.25cm);
      \draw[blue,line width=1.6pt] (0,0) circle (1.25cm);
      \node[navy,font=\bfseries\fontsize{22}{22}\selectfont] at (0,0) {CL};
    \end{tikzpicture}
  \end{minipage}%
  \hfill
  \begin{minipage}[c]{0.79\linewidth}
    {\fontsize{27}{29}\selectfont\bfseries\color{white}Camille Laurent}\par
    \vspace{1pt}
    {\Large\color{blue!25!white}Cheffe de projet digital}\par
    \vspace{7pt}
    {\small\color{white}
      \contact{map-marker-alt}{Paris, France}
      \quad
      \contact{phone-alt}{+33 6 12 34 56 78}
      \quad
      \contact{envelope}{\href{mailto:camille.laurent.cv@example.com}{\color{white}camille.laurent.cv@example.com}}
      \\[4pt]
      \contact{linkedin}{\href{https://www.linkedin.com/in/camille-laurent}{\color{white}linkedin.com/in/camille-laurent}}
      \quad
      \contact{globe-europe}{Mobilité : Île-de-France \& télétravail}
    }
  \end{minipage}
\end{tcolorbox}

\vspace{7pt}

% Corps en deux colonnes
\begin{minipage}[t]{0.305\textwidth}
  \vspace{0pt}
  \begin{tcolorbox}[
    colback=ice,
    colframe=ice,
    boxrule=0pt,
    arc=1.5mm,
    left=4.5mm,
    right=4.5mm,
    top=3mm,
    bottom=3mm,
    height=21.7cm,
    valign=top
  ]
    \RaggedRight

    \sidesection{Compétences}
    \skill{Gestion de projet}{5}
    \skill{Méthodes agiles}{5}
    \skill{Pilotage budgétaire}{4}
    \skill{Conduite du changement}{4}
    \skill{Analyse de données}{4}

    \sidesection{Outils}
    {\small
      \textbf{Produit :} Jira, Confluence, Notion\par\vspace{3pt}
      \textbf{Data :} Excel avancé, Power BI, GA4\par\vspace{3pt}
      \textbf{Collaboration :} Miro, Figma, Teams\par\vspace{3pt}
      \textbf{CRM / CMS :} Salesforce, WordPress
    }

    \sidesection{Langues}
    {\small
      \textbf{Français}\hfill Langue maternelle\par\vspace{3pt}
      \textbf{Anglais}\hfill C1 --- courant\par\vspace{3pt}
      \textbf{Espagnol}\hfill B1 --- intermédiaire
    }

    \sidesection{Certifications}
    {\small
      \textbf{Professional Scrum Master I}\par
      {\footnotesize\color{muted}Scrum.org, 2024}\par\vspace{4pt}
      \textbf{Nexa Analytics Certification}\par
      {\footnotesize\color{muted}Nexa, 2023}
    }

    \sidesection{Savoir-être}
    {\small
      Leadership collaboratif\par
      Sens de la communication\par
      Esprit d'analyse\par
      Orientation résultats
    }

    \sidesection{Centres d'intérêt}
    {\small
      \faHiking\ Randonnée itinérante\par\vspace{2pt}
      \faHandsHelping\ Mentorat associatif\par\vspace{2pt}
      \faCameraRetro\ Photographie urbaine
    }
  \end{tcolorbox}
\end{minipage}%
\hfill
\begin{minipage}[t]{0.66\textwidth}
  \vspace{0pt}
  \RaggedRight

  \mainsection{Profil}
  Cheffe de projet digital avec plus de sept ans d'expérience dans la conception
  et le déploiement de services numériques. Habituée à fédérer des équipes
  pluridisciplinaires, à traduire les besoins métier en feuilles de route
  concrètes et à piloter la performance par la donnée. Reconnue pour livrer des
  projets complexes dans les délais tout en améliorant l'expérience client.

  \mainsection{Expérience professionnelle}

  \experience{Cheffe de projet digital senior}{01/2022 -- aujourd'hui}
    {Nova Conseil}{Paris}
  \begin{itemize}
    \item Pilotage d'un portefeuille de projets e-commerce de 1,8 M\texteuro{},
      de la définition du besoin au déploiement, avec une équipe de 9 personnes.
    \item Refonte du parcours de commande ayant augmenté le taux de conversion de
      18\% et réduit les abandons sur mobile de 22\% en six mois.
    \item Mise en place d'une gouvernance agile et d'indicateurs partagés,
      diminuant de 28\% les retards de livraison entre 2022 et 2024.
  \end{itemize}

  \vspace{4pt}
  \experience{Product Owner}{09/2018 -- 12/2021}
    {Atelier Numérique}{Boulogne-Billancourt}
  \begin{itemize}
    \item Construction et priorisation du backlog de six produits internes,
      en lien avec 35 parties prenantes métiers et techniques.
    \item Automatisation du traitement des demandes clients, réduisant le délai
      moyen de réponse de 40\% et économisant 600 heures par an.
    \item Animation des ateliers utilisateurs, des revues de sprint et des comités
      de pilotage ; suivi des objectifs au moyen de tableaux de bord Power BI.
  \end{itemize}

  \vspace{4pt}
  \experience{Assistante cheffe de projet web}{03/2017 -- 08/2018}
    {Vektor Studio}{Paris}
  \begin{itemize}
    \item Coordination de la création de sites et campagnes pour 12 comptes B2B,
      avec suivi des plannings, budgets, recettes et mises en production.
    \item Formalisation d'un processus de contrôle qualité ayant réduit de 30\%
      les anomalies constatées après livraison.
  \end{itemize}

  \mainsection{Formation}

  \education{Master Management de l'innovation numérique}{2016 -- 2018}
    {Université Paris Dauphine--PSL}{Mention bien --- Paris}
  \vspace{4pt}
  \education{Licence Information et communication}{2013 -- 2016}
    {Université Paris 8 Vincennes--Saint-Denis}{Parcours médias numériques --- Saint-Denis}

  \mainsection{Projet marquant}
  \textbf{Programme d'accessibilité numérique}\hfill
  {\footnotesize\bfseries\color{muted}2023 -- 2024}\par
  {\small Définition d'une feuille de route RGAA, audit de 60 parcours et
  formation de 45 contributeurs ; passage du taux de conformité de 54\% à 82\%.}

\end{minipage}

\end{document}
